\section{Benchmark and Environment}
\label{sec:benchmark}

\subsection{Task definition}

The robot receives visual observations and must flatten a randomly crumpled
garment face-up on the table. A valid solution matches the garment's canonical
shape under one of four right-angle in-plane rotations; a back-facing garment
is not successful. Deformable distance measures shape agreement, while rigid
distance measures global alignment with the valid layouts. Their combination
provides both a success threshold and continuous progress when a rollout does
not fully solve the task.

\subsection{Initial-state distribution}

Initial states are generated by randomly orienting a garment in the air,
pinning a random point, allowing the garment to hang, and then releasing it.
This process creates folds and layer arrangements through cloth dynamics rather
than a small library of authored poses. Physically invalid states are removed,
but difficult folds, partial inversions, and self-occlusion remain. Initial
distance describes the resulting difficulty instead of acting as a filter.

\subsection{Garment distribution}

The garment set focuses on upper-body clothing but varies in sleeves, openings,
body length, scale, and dress-like structure. These variations change both what
is visible and what motions are physically useful. Semantic morphology labels
separate a new instance of a familiar form from a genuinely new structural
form. Mesh resolution is recorded but is not itself treated as task difficulty.

\subsection{Interaction environment}

The environment combines Style3D cloth simulation with continuous bimanual
robot control. The same scene and trajectory representation supports
teleoperation, policy rollout, HIL, replay, and evaluation. Dynamic operations
such as shaking and fine adjustments such as pick--place therefore coexist in
one action stream rather than being exposed as separate task primitives.

\subsection{Evaluation dimensions}

The benchmark treats initial state, garment morphology, appearance, and
dynamics as distinct sources of variation. A model can fail because it cannot
resolve a fold, cannot generalize to a garment structure, follows an appearance
shortcut, or cannot adapt its action to different dynamics. Paired conditions
reuse the same physical case when isolating one factor, while the continuous
distance metrics reveal where performance changes even when binary success is
unchanged.
